% 06-reviewers —— 由 tools/convert.py 自动生成，勿手改
% 源文件：part0000_split_017.html

作为 Rust 与系统编程领域的经典，本书翻译难度较大。在此，诚挚地感谢在第 2 版中文版出版过程中参与公开审读活动的 30 位开发者（含一线专家与爱好者）。

各位审读专家针对译文提出了大量宝贵的意见与建议，你们严谨的治学态度和对 Rust 语言的深刻理解让译者与编辑受益匪浅，也让本书的译文质量更上一层楼。在此以表格形式列出 30 位审读专家的姓名（按姓氏字母排序），各章及章名与为之贡献的专家对应。

\myTableBegin
\begin{xltabular}{\textwidth}{@{}>{\hsize=1.000\hsize\raggedright\arraybackslash}X>{\hsize=0.830\hsize\raggedright\arraybackslash}X>{\hsize=1.170\hsize\raggedright\arraybackslash}X@{}}
\rowcolor{tabhead}\textcolor{white}{审读划分} & \textcolor{white}{章名} & \textcolor{white}{审读专家} \\
\endfirsthead
\rowcolor{tabhead}\textcolor{white}{审读划分} & \textcolor{white}{章名} & \textcolor{white}{审读专家} \\
\endhead
\multirow{3}{106.0bp}{第 1\textasciitilde{}3 章} & 第 1 章 系统程序员也能享受美好 & \multirow{3}{126.4bp}{陈骜、姜子龙、尚卓燃（PsiACE）、唐刚（@Mike Tang）} \\
 & 第 2 章  Rust 导览 &  \\
 & 第 3 章 基本数据类型 &  \\
\multirow{3}{106.0bp}{第 4\textasciitilde{}6 章} & 第 4 章 所有权与移动 & \multirow{3}{126.4bp}{陈翔（@翔翔的学习频道）、刘燚、陶叶港（@scruel）、杨思杰} \\
 & 第 5 章 引用 &  \\
 & 第 6 章 表达式 &  \\
\multirow{3}{106.0bp}{第 7\textasciitilde{}9 章} & 第 7 章 错误处理 & \multirow{3}{126.4bp}{寸志、胡屹、苏胤榕、张卫滨} \\
 & 第 8 章  crate 与模块 &  \\
 & 第 9 章 结构体 &  \\
\multirow{3}{106.0bp}{第 10\textasciitilde{}12 章} & 第 10 章 枚举与模式 & \multirow{3}{126.4bp}{刘久武、张鑫明、阴钰（@yxonic）} \\
 & 第 11 章 特型与泛型 &  \\
 & 第 12 章 运算符重载 &  \\
\multirow{3}{106.0bp}{第 13\textasciitilde{}15 章} & 第 13 章 实用工具特型 & \multirow{3}{126.4bp}{曹洪伟、柳佳龙、尹钢、张罗东（@zhangll）} \\
 & 第 14 章 闭包 &  \\
 & 第 15 章 迭代器 &  \\
\multirow{3}{106.0bp}{第 16\textasciitilde{}18 章} & 第 16 章 集合 & \multirow{3}{126.4bp}{Kyle C、杜小豆、钱宇超（@stdrc）、杨光宇} \\
 & 第 17 章 字符串与文本 &  \\
 & 第 18 章 输入与输出 &  \\
\multirow{2}{106.0bp}{第 19\textasciitilde{}20 章} & 第 19 章 并发 & \multirow{2}{126.4bp}{耿腾、刘祺、逄振洲、赵梓淇（BugenZ）} \\
 & 第 20 章 异步编程 &  \\
\multirow{3}{106.0bp}{第 21\textasciitilde{}23 章} & 第 21 章 宏 & \multirow{3}{126.4bp}{艾尼、大妈（ZoomQuiet）、刘鑫} \\
 & 第 22 章 不安全代码 &  \\
 & 第 23 章 外部函数 &  \\
\end{xltabular}
\myTableEnd

本书虽已出版，但内容品质的提升不会终止。译者、编辑、审读专家虽已尽力，但疏漏可能在所难免，如果大家在阅读过程中发现任何问题，欢迎将其提交到图灵社区本书的勘误处（ituring.com.cn/\allowbreak{}book/\allowbreak{}2846）。勘误经编辑确认之后会更正在重印书中。

